\setchapterpreamble[u]{\margintoc}
\chapter{Roteiro de Estudo}
\labch{cet615-roteiro}

Perguntas para revisar a disciplina, capítulo por capítulo. As respostas
não estão aqui: o objetivo é voltar ao texto e às referências para
encontrá-las. As perguntas têm três níveis.

\begin{description}
	\item[Rever.] A resposta está no capítulo indicado. Serve para conferir
	se o conceito ficou.
	\item[Aplicar.] Pede para escrever, corrigir ou comparar modelos e
	comandos num cenário novo, com o que o capítulo ensina. O SQL Server é
	o melhor juiz: escreva, execute e veja o resultado.
	\item[Aprofundar.] Vai além da aula. A resposta exige as referências
	citadas no capítulo, o \textcite{Elmasri2011} ou a documentação do SQL
	Server, e muitas vezes admite mais de uma resposta defensável, desde
	que justificada.
\end{description}

A sugestão é responder por escrito antes de conferir. Uma resposta que
não cabe em poucas linhas costuma indicar que o conceito ainda não
assentou.

\section{Conceitos fundamentais}

Referente ao \refch{cet615-conceitos}.

\subsection*{Rever}
\begin{enumerate}
	\item O que distingue um banco de dados de uma coleção qualquer de
	dados? O que é o minimundo?
	\item Quais problemas do sistema de arquivos o SGBD resolve?
	\item Quais são as fases do projeto de um banco de dados, e o que cada
	uma entrega para a seguinte?
	\item Em que dois casos a aula recomenda usar generalização?
\end{enumerate}

\subsection*{Aplicar}
\begin{enumerate}
	\item Modele em ME-R uma clínica com médicos, pacientes e consultas.
	Indique grau, cardinalidade e chaves. Onde caberia uma generalização?
	\item Mapeie para o modelo relacional a generalização de pessoa em
	aluno e professor de duas formas diferentes. Compare as tabelas que
	resultam de cada uma.
	\item Dê um exemplo de linha que viola uma chave primária e outro de
	linha que viola uma chave estrangeira no banco da biblioteca.
\end{enumerate}

\subsection*{Aprofundar}
\begin{enumerate}
	\item Pela arquitetura de três esquemas, qual das duas independências
	de dados é mais difícil de alcançar na prática, e por quê?
	\item Entre as opções de mapeamento de especialização do Elmasri,
	quando uma tabela única com atributo de tipo é melhor que uma tabela
	por específico? Considere nulos, junções e restrições.
\end{enumerate}

\section{O SQL Server e o banco de dados}

Referente ao \refch{cet615-sqlserver}.

\subsection*{Rever}
\begin{enumerate}
	\item Para que servem os bancos \texttt{master}, \texttt{model},
	\texttt{msdb} e \texttt{tempdb}?
	\item Quais arquivos compõem um banco no SQL Server, e o que cada um
	guarda?
	\item O que o \texttt{go} faz, e por que ele não é um comando T-SQL?
\end{enumerate}

\subsection*{Aplicar}
\begin{enumerate}
	\item Escreva um \texttt{create database} com um arquivo de dados e um
	de log, com tamanho inicial e crescimento definidos. Depois aumente o
	arquivo de dados com \texttt{alter database}.
	\item Crie uma tabela no \texttt{model} e depois um banco novo. Onde a
	tabela aparece? Explique o resultado.
	\item Use as visões de catálogo para listar as tabelas e as colunas de
	um banco seu.
\end{enumerate}

\subsection*{Aprofundar}
\begin{enumerate}
	\item Que consequências o fato de o \texttt{tempdb} ser recriado a cada
	reinício tem para quem usa tabelas temporárias e para o desempenho da
	instância?
	\item Quando vale distribuir um banco em mais de um grupo de arquivos?
	Pense em discos diferentes, backup e partes do banco só de leitura.
\end{enumerate}

\section{Definição de dados (DDL)}

Referente ao \refch{cet615-ddl}.

\subsection*{Rever}
\begin{enumerate}
	\item Por que basta uma coluna \texttt{varchar} para o registro passar
	a ter tamanho variável? Quando a aula recomenda \texttt{varchar}?
	\item Quais restrições de integridade podem ser declaradas no
	\texttt{create table}?
	\item Por que só pode haver um índice clusterizado por tabela?
\end{enumerate}

\subsection*{Aplicar}
\begin{enumerate}
	\item Escreva o script de criação das tabelas livro, leitor e
	empréstimo, com chaves, restrições nomeadas e um \texttt{check}.
	\item Proponha três consultas frequentes sobre essas tabelas e escolha
	a partir delas o índice clusterizado e os não clusterizados.
	Justifique cada escolha.
	\item Remova uma chave estrangeira nomeada, grave uma linha que a
	viola e tente recriar a chave. O que acontece?
\end{enumerate}

\subsection*{Aprofundar}
\begin{enumerate}
	\item Para a chave estrangeira do empréstimo para o leitor, compare
	\texttt{cascade}, \texttt{set null} e a ausência de ação referencial.
	Qual você escolheria, e por quê?
	\item Um índice composto em (\texttt{sobrenome}, \texttt{nome}) ajuda
	numa consulta que filtra só por \texttt{nome}? Explique com a
	estrutura da árvore B.
\end{enumerate}

\section{Manipulação de dados (DML)}

Referente ao \refch{cet615-dml}.

\subsection*{Rever}
\begin{enumerate}
	\item Qual a diferença entre um \texttt{delete} sem \texttt{where} e um
	\texttt{truncate table}?
	\item Como a lógica de três valores afeta um \texttt{where} sobre uma
	coluna com nulos?
	\item Quais são os estados de uma transação e as propriedades ACID?
	\item Qual erro de junção a aula diz que conta na nota, mesmo com o
	resultado certo?
\end{enumerate}

\subsection*{Aplicar}
\begin{enumerate}
	\item Liste os leitores que nunca pegaram um livro emprestado de duas
	formas: com junção externa e com subconsulta.
	\item Para cada autor com mais de dois livros, mostre a quantidade de
	livros. Na sua consulta, o que vai no \texttt{where} e o que vai no
	\texttt{having}?
	\item Em duas janelas do SSMS, abra uma transação que altera um
	empréstimo e, sem confirmar, consulte a mesma linha na outra janela.
	O que acontece, e por quê?
\end{enumerate}

\subsection*{Aprofundar}
\begin{enumerate}
	\item Por que o log de transações precisa ser gravado antes das
	páginas de dados? O que isso permite na recuperação depois de uma
	queda?
	\item O comportamento observado no exercício das duas janelas depende
	do nível de isolamento. Quais níveis o SQL Server oferece, e o que
	cada um troca entre consistência e concorrência?
\end{enumerate}

\section{Visões}

Referente ao \refch{cet615-visoes}.

\subsection*{Rever}
\begin{enumerate}
	\item O que é uma visão? Quais as vantagens e as desvantagens?
	\item Em que condições uma visão é atualizável?
	\item Por que a professora conclui que visão serve para consulta?
\end{enumerate}

\subsection*{Aplicar}
\begin{enumerate}
	\item Crie uma visão dos empréstimos em atraso com o nome do leitor e
	o título do livro, e uma consulta a ela que filtre por leitor.
	\item Uma colega usa essa visão com \texttt{select distinct} para
	listar os leitores com empréstimo. Qual o problema, e qual a consulta
	certa?
	\item Tente inserir uma linha por uma visão que omite uma coluna
	\texttt{not null}. Explique o erro.
\end{enumerate}

\subsection*{Aprofundar}
\begin{enumerate}
	\item Que problema o \texttt{with check option} resolve? Mostre um
	caso em que, sem ele, a linha alterada some da própria visão.
	\item O que muda quando uma visão passa a ser indexada? Quando o
	custo de manter o índice compensa?
\end{enumerate}

\section{Procedimentos armazenados e transações}

Referente ao \refch{cet615-procedimentos}.

\subsection*{Rever}
\begin{enumerate}
	\item Quais as vantagens de um procedimento armazenado?
	\item Qual a diferença entre parâmetro de saída e valor de retorno?
	\item Por que não dá para escrever um único \texttt{return
	@@rowcount} no fim do procedimento?
\end{enumerate}

\subsection*{Aplicar}
\begin{enumerate}
	\item Escreva um procedimento que registre um empréstimo e atualize a
	situação do exemplar numa única transação, retornando 1 ou 0. Teste
	um caso de sucesso e dois de falha.
	\item Acrescente a esse procedimento um parâmetro com valor padrão e
	um de saída. Mostre uma chamada que use os dois.
	\item Um procedimento chama outro, e o interno executa
	\texttt{rollback}. O que acontece com o \texttt{@@trancount} e com o
	procedimento externo?
\end{enumerate}

\subsection*{Aprofundar}
\begin{enumerate}
	\item Compare o padrão com \texttt{@@rowcount} usado em aula com
	\texttt{try...catch} e \texttt{throw}. Que erros cada um deixa
	escapar?
	\item Por que o padrão SQL e o T-SQL tratam procedimentos de forma tão
	diferente? Que consequência isso tem para levar um sistema de um SGBD
	para outro?
\end{enumerate}

\section{Gatilhos}

Referente ao \refch{cet615-gatilhos}.

\subsection*{Rever}
\begin{enumerate}
	\item O que é um banco de dados ativo? O que significa cada parte do
	modelo evento--condição--ação?
	\item O que as tabelas \texttt{inserted} e \texttt{deleted} contêm em
	um \texttt{insert}, um \texttt{update} e um \texttt{delete}?
	\item Por que um \texttt{rollback} dentro do gatilho desfaz também o
	comando que o disparou?
	\item Qual teste a professora propôs para decidir entre gatilho e
	transação explícita?
\end{enumerate}

\subsection*{Aplicar}
\begin{enumerate}
	\item Escreva os gatilhos que mantêm a quantidade de empréstimos
	ativos de cada leitor, preparados para comandos que afetam várias
	linhas.
	\item O gatilho da segunda avaliação falha num \texttt{insert} com duas
	compras. Reescreva-o para várias linhas e teste.
	\item Com a chave estrangeira do item para a nota, o gatilho que apaga
	os itens de uma nota excluída nunca roda. Demonstre e explique.
\end{enumerate}

\subsection*{Aprofundar}
\begin{enumerate}
	\item Em que situação um gatilho \texttt{instead of} é a única saída,
	e não apenas um jeito de simular \texttt{before}? Pense em visões.
	\item Quais os riscos de gatilhos encadeados e recursivos, e como o
	SQL Server os limita?
\end{enumerate}

\section{Controle de dados (DCL)}

Referente ao \refch{cet615-dcl}.

\subsection*{Rever}
\begin{enumerate}
	\item Quais são os três níveis de acesso, e o que acontece quando
	falta um deles?
	\item Qual a diferença entre \texttt{grant}, \texttt{revoke} e
	\texttt{deny}? Qual tem precedência?
	\item Quais são os tipos de backup, e o que cada um copia?
\end{enumerate}

\subsection*{Aplicar}
\begin{enumerate}
	\item Crie um login, um usuário e um papel de atendentes que leia
	livros e registre empréstimos. Negue a exclusão de empréstimos a um
	atendente específico e teste conectado como ele.
	\item Monte uma estratégia de backup para um banco que não pode perder
	mais de uma hora de dados. Quais tipos, com que frequência, e qual a
	sequência de restauração depois de uma falha?
	\item Use um esquema para separar os objetos do empréstimo e conceda
	permissão no esquema inteiro, em vez de objeto por objeto.
\end{enumerate}

\subsection*{Aprofundar}
\begin{enumerate}
	\item Em que casos o \texttt{deny} não vence? Por que essas exceções
	existem?
	\item Como o modelo de recuperação (\emph{simple}, \emph{full} ou
	\emph{bulk-logged}) limita a estratégia de backup que você montou?
\end{enumerate}
